\uuid{Zy4V}
\titre{Identification de graphes de fonctions}
\niveau{L2}
\module{Analyse}
\chapitre{Fonctions de plusieurs variables}
\sousChapitre{Représentation graphique}
\theme{Graphes de fonctions de deux variables, reconnaissance visuelle}
\auteur{Maxime Nguyen}
\datecreate{2026-03-24}
\organisation{amscc}
\difficulte{2}

\contenu{
\texte{
  Un élève a collecté les graphes de fonctions des exercices précédents mais ne se souvient plus à quel exercice correspond chaque graphe. Pour chacune des figures suivantes, donner le numéro de l'exercice contenant la fonction correspondante.
}
\begin{enumerate}
  \item \question{Figure 1.}
    \reponse{Figure 1 : exercice 5. On reconnaît le domaine de définition et la courbe de niveau $0$ de $f(x,y) = \dfrac{\sqrt{5y-2x}}{\ln(x-1)}$.}

  \item \question{Figure 2.}
    \reponse{Figure 2 : exercice 4. On reconnaît les points critiques formant l'axe $\{(0,y) \mid y \in \mathbb{R}\}$ de $f(x,y) = x^2 y$.}

  \item \question{Figure 3.}
    \reponse{Figure 3 : exercice 8. On reconnaît le maximum ainsi que la courbe de niveau $0$ de la fonction d'efficacité $E(x,y)$.}

  \item \question{Figure 4.}
    \reponse{Figure 4 : exercice 12. On reconnaît le minimum global de $f(x,y) = 9x^2 + 8xy + 3y^2 + x - 2y$.}
\end{enumerate}
}
